\begingroup
\renewcommand{\arraystretch}{1.45}
\fancyhead[L]{\small ARITHMETISCH-GEOMETRISCHES MITTEL.}
\Needspace{28\baselineskip}
\AGMArticle{21}

[Bildet man die Gleichungen zwischen den Ableitungen von $\alpha,\allowbreak \beta,\allowbreak \gamma,\allowbreak \delta$ nach der Größe $y$ als unabhängiger Variablen, so müssen dieselben Koeffizienten dieser Ableitungen entstehen wie in den Gleichungen für die Differentiale nach $\eta$. Da sich die letzteren Gleichungen so darstellen lassen, dass die Verhältnisse zwischen den Ableitungen von Quotienten bestimmt werden, müssen die genannten Koeffizienten auch zu diesen Ableitungen umgekehrt proportional sein. Ersetzt man sie durch deren Kehrwerte, so erkennt man, dass auch die übrigen Glieder jener Gleichungen durch solche Ableitungen dargestellt werden können, und zwar so, dass der Ausdruck von der Form
\Needspace{10\baselineskip}
\[
\dfrac{\left[\dfrac{\mathrm{d}\log\dfrac{\delta}{\beta}}{\mathrm{d}\log y}\right]}
{\left[\dfrac{\mathrm{d}\log\dfrac{\delta}{\beta}}{\mathrm{d}\log\eta}\right]}
-\tfrac12\dfrac{\mathrm{d}\log\dfrac{\alpha\beta}{xx}}{\mathrm{d}\log\eta}
\]
für einen beliebigen gemeinsamen Index der Zeichen $\alpha,\allowbreak \beta,\allowbreak \gamma,\allowbreak \delta,\allowbreak x,\allowbreak y,\allowbreak \eta$ seinen Wert nicht ändert, in welcher Reihenfolge die Größen $\alpha,\allowbreak \beta,\allowbreak \gamma,\allowbreak \delta$ auch miteinander vertauscht werden. Vergleicht man diesen Ausdruck, der sich auf $\frac{\delta}{\beta},\allowbreak \frac{\alpha\gamma}{xx},\allowbreak \eta,\allowbreak y$ bezieht, mit demjenigen für $\frac{\delta'}{\beta'}=\frac{\gamma\delta}{\alpha\beta},\allowbreak \frac{\alpha'\gamma'}{x'x'},\allowbreak \eta'=\eta\eta,\allowbreak y'=yy$ und beachtet, dass nach Artikel 19
\Needspace{14\baselineskip}
\[
\begin{gathered}\left[\dfrac{\mathrm{d}\log\dfrac{\alpha}{\beta}}{\mathrm{d}\log\eta}\right]^2
=\dfrac{\mathrm{d}\log\dfrac{\gamma'}{\beta'}}{\mathrm{d}\log\eta}\cdot
\dfrac{\mathrm{d}\log\dfrac{\alpha'}{\beta'}}{\mathrm{d}\log\eta},\\[9pt]
\left[\dfrac{\mathrm{d}\log\dfrac{\delta}{\gamma}}{\mathrm{d}\log\eta}\right]^2
=\dfrac{\mathrm{d}\log\dfrac{\delta'}{\alpha'}}{\mathrm{d}\log\eta}\cdot
\dfrac{\mathrm{d}\log\dfrac{\delta'}{\gamma'}}{\mathrm{d}\log\eta}\end{gathered}
\]
gilt, so ergibt sich, dass der Wert jenes Ausdrucks auch vom Index der Zeichen $\alpha,\allowbreak \beta,\allowbreak \gamma,\allowbreak \delta,\allowbreak x,\allowbreak y,\allowbreak \eta$ unabhängig und, wie aus den Grenzwerten folgt, gleich null sein muss.
\Needspace{22\baselineskip}
\SourceLeaf{402}

Um die Ableitung nach $y$ jeder einzelnen der vier Größen $\alpha,\allowbreak \beta,\allowbreak \gamma,\allowbreak \delta$ zu bestimmen, wird es zunächst erforderlich sein, jenen Ausdruck so umzuformen, dass er nur zwei der Größen enthält. Dies kann zum Beispiel mithilfe der im vorhergehenden Artikel bestimmten, nach $\eta$ genommenen zweiten Ableitungen der Quotienten geschehen, sodass alle Ausdrücke
\Needspace{12\baselineskip}
\[
\begin{aligned}&\dfrac{\mathrm d\log\dfrac{\alpha}{\beta}}{\mathrm d\log y}
-\dfrac{\mathrm d\mathrm d\log\dfrac{\alpha}{\beta}}{(\mathrm d\log\eta)^2}
\\[8pt]&{}-\dfrac12\dfrac{\mathrm d\log\dfrac{\alpha\beta}{xx}}{\mathrm d\log\eta}
\cdot\dfrac{\mathrm d\log\dfrac{\alpha}{\beta}}{\mathrm d\log\eta}\end{aligned}
\]
verschwinden, in welcher Reihenfolge $\alpha,\allowbreak \beta,\allowbreak \gamma,\allowbreak \delta$ auch miteinander vertauscht werden und welcher gemeinsame Index ihnen und den Zeichen $y,\allowbreak x,\allowbreak \eta$ gegeben wird. Ersetzt man $\frac{\alpha}{\beta}$ durch $\frac{b'\beta'}{\beta\beta}$, $\alpha\beta$ durch $b'\beta'$ und nach dem vorhergehenden Artikel
\Needspace{14\baselineskip}
\[
\begin{gathered}\dfrac{\mathrm d\log\dfrac{b'}{k}}{2\mathrm d\log y}
\quad\text{durch}\\[9pt]
\dfrac{\mathrm d\mathrm d\log\dfrac{\beta'}{x}}{(2\mathrm d\log\eta)^2}
+\dfrac12\left[\dfrac{\mathrm d\log\dfrac{\beta\beta}{\beta'}}{\mathrm d\log\eta}\right]^2\end{gathered}
\]
so erkennt man, dass der Ausdruck
\Needspace{12\baselineskip}
\[
\begin{aligned}&\dfrac{\mathrm d\log\dfrac{\beta}{x}}{\mathrm d\log y}
-\dfrac{\mathrm d\mathrm d\log\dfrac{\beta}{x}}{(\mathrm d\log\eta)^2}
\\[8pt]&{}-\dfrac12\left[\dfrac{\mathrm d\log\dfrac{\beta}{x}}{\mathrm d\log\eta}\right]^2\end{aligned}
\]
seinem Wert nach unverändert bleibt, nicht nur, wie eben gefunden, wenn $\beta$ mit $\alpha$ oder $\gamma$ oder $\delta$ vertauscht wird, sondern auch, wenn diesen Größen zugleich mit $y,\allowbreak x,\allowbreak \eta$ ein beliebiger gemeinsamer Index gegeben wird. Für einen beständig wachsenden Index verschwindet der Wert des Ausdrucks, und durch Multiplikation mit $\frac12\surd {\frac{\beta}{x}}$ entsteht demnach
\Needspace{10\baselineskip}
\[
\dfrac{\mathrm d\surd {\dfrac{\alpha}{x}}}{\mathrm d\log y}
-\dfrac{\mathrm d\mathrm d\surd {\dfrac{\alpha}{x}}}{(\mathrm d\log\eta)^2}=0
\]
als partielle Differentialgleichung für jede der Größen $\alpha,\allowbreak \beta,\allowbreak \gamma,\allowbreak \delta$.]

\Needspace{26\baselineskip}
\AGMArticle{22}
[Die Werte aller Glieder des Algorithmus hängen von vier Größen ab. Als solche dürfen $a,\allowbreak b,\allowbreak \alpha,\allowbreak \beta$ angenommen werden, aber auch $k,\allowbreak y,\allowbreak x,\allowbreak \eta$, weil diese sich unabhängig voneinander ändern können. Sämtliche Gleichungen zwischen den Gliedern der Folgen lassen sich so darstellen, dass sie sowohl in Bezug auf $a,\allowbreak b,\allowbreak c,\allowbreak .\;. a^{\mathrm n},\allowbreak b^{\mathrm n},\allowbreak c^{\mathrm n},\allowbreak k$ als auch in Bezug auf


\SourceLeaf{403}

\Needspace{10\baselineskip}
\[
\begin{gathered}\dfrac{\alpha}{k},\ \dfrac{\beta}{k},\ \dfrac{\gamma}{k},\ \dfrac{\delta}{k}\ .\;.\;.\\[8pt]\
{}^{2^n}\!\surd {\dfrac{\alpha^{\mathrm n}}{k}},\
{}^{2^n}\!\surd {\dfrac{\beta^{\mathrm n}}{k}},\
{}^{2^n}\!\surd {\dfrac{\gamma^{\mathrm n}}{k}},\
{}^{2^n}\!\surd {\dfrac{\delta^{\mathrm n}}{k}},\ \dfrac{x}{k}\end{gathered}
\]
homogen sind. Die Verhältnisse zwischen den letzteren hängen also nicht von $k$ und $x$, sondern allein von $y$ und $\eta$ ab; wir können also
\Needspace{14\baselineskip}
\[
\begin{gathered}\alpha=x P(y,\eta)^2,\\[5pt]
\beta=x Q(y,\eta)^2,\\[5pt]
\gamma=x R(y,\eta)^2,\\[5pt]
\delta=x S(y,\eta)^2\end{gathered}
\]
setzen. Aus den Betrachtungen in Artikel 18 folgt dann, dass diese Funktionen $P,\allowbreak \ Q,\allowbreak \ R,\allowbreak \ S$ dieselben bleiben, wenn $\alpha,\allowbreak \beta,\allowbreak \gamma,\allowbreak \delta,\allowbreak x,\allowbreak y,\allowbreak \eta$ mit einem beliebigen Index versehen werden. \Needspace{16\baselineskip}
Aus demselben Artikel folgt auch, mit der in Artikel 15 verwendeten Bezeichnung, für die mit einem gleichen Index versehenen $a,\allowbreak b. c,\allowbreak y$
\[
\begin{aligned}
P(y,1)&=py=\surd {\dfrac ak}=\surd {\dfrac a{\mathrm M(a,b)}}\\
Q(y,1)&=qy=\surd {\dfrac bk}\\
R(y,1)&=ry=\surd {\dfrac ck}\\
S(y,1)&=0
\end{aligned}
\]
ferner als Gleichungen, die denjenigen entsprechen, durch die die Größen $\gamma$ und $\delta$ bestimmt wurden:
\Needspace{24\baselineskip}
\[
\begin{gathered}
qqrr(qqRR-rrQQ)\\[4pt]=rrpp(ppRR-rrPP)\\
=ppqq(qqPP-ppQQ)\\[4pt]=ppqqrrSS\\[4pt]
\dfrac{PP+QQ}{p(yy)}=\dfrac{RR-SS}{r(yy)}.\\[9pt]
\dfrac{PP-QQ}{r(yy)}=\dfrac{RR+SS}{p(yy)}\\[4pt]
\dfrac{PP+RR}{p(\surd  y)}=\dfrac{QQ+SS}{q(\surd  y)},\\[9pt]
\dfrac{PP-RR}{q(\surd  y)}=\dfrac{QQ-SS}{p(\surd  y)}
\end{gathered}
\]
\Needspace{12\baselineskip}
und als Gleichungen, die das Bildungsgesetz des Algorithmus darstellen:
\[
\begin{aligned}
2p(yy)\cdot P(yy,\eta\eta)&=PP+QQ,\\[5pt]
2p(yy)\cdot R(yy,\eta\eta)&=RR+SS\\[5pt]
2r(yy)\cdot R(yy,\eta\eta)&=PP-QQ,\\[5pt]
2r(yy)\cdot P(yy,\eta\eta)&=RR-SS\\[5pt]
q(yy)\cdot Q(yy,\eta\eta)&=PQ,\\[5pt]
q(yy)\cdot S(yy,\eta\eta)&=RS
\end{aligned}
\]
wobei alle Funktionen, neben denen kein Argument steht, sich auf die einfachen $y$ und $\eta$ beziehen.

\Needspace{22\baselineskip}
\SourceLeaf{404}

Zur Berechnung der Werte der Funktionen, die zu gegebenen Werten von $\frac ab=\frac{pp}{qq}$ und $U$ oder $V$ gehören, erhält man aus der Formel von \textsc{Gauss} für $H$ oder $x$, wenn man zur Abkürzung
\Needspace{10\baselineskip}
\[
\begin{aligned}&(a^{\mathrm n}\sin 2^{\mathrm n} U^{\mathrm n})^2+(b^{\mathrm n}\cos 2^{\mathrm n} U^{\mathrm n})^2\\[5pt]&=b^{\mathrm n} b^{\mathrm n}\Delta^{\mathrm n}\Delta^{\mathrm n}\end{aligned}
\]
setzt, die Gleichungen:
\Needspace{18\baselineskip}
\[
\begin{aligned}
Q&=q\cdot\surd \Delta\cdot{}^{4}\!\surd {\Delta'}\cdot{}^{8}\!\surd {\Delta''}\cdot{}^{16}\!\surd {\Delta'''}.\;.\;.\\[9pt]
P&=\dfrac pq Q\dfrac1\Delta\\[9pt]
R&=\dfrac rq Q\dfrac{\cos U}{\Delta}\\[9pt]
S&=\dfrac pr Q\dfrac{\sin U}{\Delta}i
\end{aligned}
\]

\Needspace{20\baselineskip}
Die Bestimmung der Größen $k,\allowbreak y,\allowbreak x,\allowbreak \eta$ als Grenzwerte zeigt, dass bei geeigneter, hier noch zulässiger Wahl der Vorzeichen der Funktionen $P,\allowbreak Q,\allowbreak R,\allowbreak S$ für bis auf null abnehmende Werte von $y,\allowbreak y\eta,\allowbreak \frac1\eta$ die Ausdrücke
\Needspace{16\baselineskip}
\[
\begin{gathered}P(y,\eta),\\[6pt] Q(y,\eta).\\[6pt]
\dfrac{R(y,\eta)}{y^{\frac14}\eta^{\frac12}},\\[6pt]
\dfrac{S(y,\eta)}{y^{\frac14}\eta^{\frac12}}\end{gathered}
\]
\Needspace{16\baselineskip}
sich dem Grenzwert eins nähern, dass aber für ein beständig abnehmendes $y$ und ein endliches reelles $u$ die Ausdrücke
\[
\begin{gathered}P(y,e^{2iu}),\\[6pt] Q(y,e^{2iu}),\\[6pt]
\dfrac{R(y,e^{2iu})}{2y^{\frac14}\cos u},\\[6pt]
\dfrac{S(y,e^{2iu})}{i2y^{\frac14}\sin u}\end{gathered}
\]
jenen Grenzwert haben.

\Needspace{16\baselineskip}
Die Funktionen $P,\allowbreak Q,\allowbreak R,\allowbreak iS$ haben reelle Werte für ein komplexes $\eta^{\frac12}$ von der Form $e^{iu}$ und bleiben bis auf $iS$, das nur sein Vorzeichen wechselt, unverändert, wenn man $u$ in $-u$ umwandelt; es gilt also
\Needspace{10\baselineskip}
\[
\begin{aligned}\dfrac{P(y,\dfrac1\eta)}P
&=\dfrac{Q(y,\dfrac1\eta)}Q
\\[7pt]&=\dfrac{R(y,\dfrac1\eta)}R
\\[7pt]&=\dfrac{-S(y,\dfrac1\eta)}S=1
\text{]}\end{aligned}
\]

\clearpage
\endgroup
